\section{Introduction}

LLZO research links a nominal garnet composition to several distinct questions: how a precursor mixture becomes a powder, how composition and substitution relate to phase state, how a powder becomes a ceramic, and how the resulting specimen is measured. A broad review in the audited set likewise spans crystal structure, preparation, doping, and conductivity rather than treating synthesis as an isolated operation \cite{raju2021crystal}. The difficulty is therefore not a shortage of processing labels. It is that a label such as ``solid state,'' ``sol--gel,'' or ``spark plasma sintering'' rarely specifies the complete chain that determines what was ultimately compared.

The stored abstracts illustrate this coupling without licensing a universal mechanism. A solid-state processing study reports sensitivity to lithium inventory and powder geometry along its route to Al-containing LLZO \cite{heywood2023tailoring}. A solution-based study focuses on precursor homogeneity and follows compositional and structural evolution \cite{parascos2022compositional}. A densification study brings powder-size distribution, heating schedule, support geometry, spatial uniformity, and transport into the same reported case \cite{guo2023uniform}. These examples make a route-only catalogue inadequate: upstream composition, firing boundary conditions, specimen geometry, and measurement definitions can all change at once.

This review therefore contributes an A+B+C bundle. Contribution A is a faceted process framework whose primary chain is \emph{powder synthesis $\rightarrow$ phase\slash composition control $\rightarrow$ thermal treatment\slash densification $\rightarrow$ structure\slash transport readout}. Composition and environment variables remain cross-cutting modifiers. Contribution B is a minimum comparability record that keeps nominal and measured composition, precursors, thermal history, environment, structure, transport definition, and replication in one account. Contribution C is a problem map in which each gap is written as an observed pattern, an evidence limitation, and a falsifiable next action.

The aim is interpretive discipline, not a performance leaderboard. We ask what the present collection permits a reader to associate, what it does not permit a reader to rank, and which missing fields prevent direct comparison. Figure~\ref{fig:process-map} provides the navigation map; Table~\ref{tab:min-record} provides the comparison rules. The scope is restricted to LLZO\slash LLZTO synthesis, phase stabilization, densification, process environment, and structure--transport evidence. Interface or cell observations appear only where their source record also reports processing-relevant LLZO information.
